% ==== partapx 附录 A:Riemann 与 Lorentz 几何的基本概念 + 参考文献 ====
% 源: MinerU_latex_A-非.tex 第 718-2283 行;校对 pdf_text.txt 第 11552-12658 行(PDF 第 213-234 页)
% 术语: 时空=spacetime, 符号差=signature, 可定向的=orientable, 切空间=tangent space,
%   惯性坐标=inertial coordinates, 零锥=null cone, 类光的=null, 零对=null pair,
%   外微分=exterior derivative, Lie 括号=Lie bracket, k 维分布=k-distribution,
%   对合的=involutive, 可积的=integrable, Levi-Civita 联络=Levi-Civita connection,
%   Christoffel 符号=Christoffel symbols, 平行移动=parallel transport, 仿射参数=affine parameter,
%   世界线=world line, 固有时=proper time, 指数映射=exponential map, 测地完备的=geodesically complete,
%   法邻域=normal neighbourhood, Weyl 张量=Weyl tensor, 共形的=conformal, 等距=isometry,
%   共形等距=conformal isometry, 共形 Killing 向量场=conformal Killing vector field,
%   变形张量=deformation tensor, 体积形式=volume form, 散度定理=divergence theorem,
%   透镜形区域=lens shaped domain, D'Alembert 算子=D'Alembertian, 时间函数=time function,
%   程函方程=eikonal equation, 因果未来=causal future, 因果过去=causal past,
%   保守区域=conservative domain, 零测地线生成元=null geodesic generator,
%   零基本形式=null fundamental form, 测地叶状结构=geodesic foliation,
%   膨胀系数=expansion coefficient, 零时距函数=null lapse function,
%   正能量条件=positive energy condition, 无散度的=divergenceless, 焦散=caustic,
%   逆平移=inverted translations, 线汇=congruence, Poincaré 群=Poincaré group,
%   共形群=conformal group, 共形不变的=conformally invariant
% REFS:
%   - 源 986 行 "See [?], Hawking and Ellis, section 2.4" 与源 1347-1348 行
%     "(see Hawking-Ellis, [?], chapter 2, section 2.6)": 原书 [?] 引文占位符缺失,已省略
%   - 源 1392 行 "citeHawkEll, chapter 2, section 2.6": OCR 破损引用,按 Hawking--Ellis 第 2 章第 2.6 节处理
%   - 源 2207 行 "According to eqrefeq:cons-law1": 破损引用,按上下文解析为 \eqref{eq:cA-443}(即刚推出的 (443))
%   - 脚注 1 中 "The large scale structure of space-time, [?]" 的 [?] 占位符省略
%   - OCR-BROKEN: 源 1439 行 D/dt 系统;源 1595 行分母 (k-1)!(原书笔误,应作 (n-1)!,照抄);
%     源 1779 行 Σ_2 也取 t_1(原书笔误,应作 t_2,照抄);源 1652 行 "12lodtsn" 系 OCR 乱码,
%     还原为 ε_{12...n}
% 编号说明:
%   - 公式 (401)-(446) 共 46 个,按原文顺序自动编号;原书未编号的行间公式用 [ ] 不占号
%   - 定理类编号与原书一致: 第 2 节 2.1-2.6、第 3 节 3.1-3.8、第 4 节 4.1-4.7
%   - 原书第 3 节 3.2 内的三个未编号 Definition 与一个未编号 Remark 用带 * 环境
%   - 原书第 4 节 Theorem 4.3 之后的未编号 Remark(437 式)用 remark*
%   - 附录公式从 (401) 起自动接续(主文件前 9 章公式止于 (400))
%   - 参考文献 [Kl-Rodn2] 两条显示标签相同,照原书;内部键分 kl-rodn2a / kl-rodn2b
% ====
\chapter{Riemann 与 Lorentz 几何的基本概念}

\section{引言}

在下面的内容中, 我们简要概述 Riemann 与 Lorentz 几何的基本概念. 关于更详细的综述, 我们分别参见 \emph{\cite{pet}}(Riemann 几何)以及 \emph{\cite{car}}、\emph{\cite{ha-e}}、\emph{\cite{wa}}(Lorentz 几何).

一个\newterm{pseudo-riemannian-manifold}{伪 Riemann 流形}{pseudo-Riemannian manifold}\footnote{我们假定读者已经熟悉微分几何的基本概念, 如流形、张量场、协变、Lie 与外微分. 关于这些概念的简短介绍, 参见 Hawking 与 Ellis 的《The large scale structure of space-time》第 2 章.}, 或简称为时空, 由一对 $(\mathbf{M}, \mathbf{g})$ 组成, 其中 $\mathbf{M}$ 是一个\newterm{orientable}{可定向的}{orientable} $p+q=n$ 维流形, $\mathbf{g}$ 是定义在其上的伪 Riemann 度量, 即一个光滑的、非退化的、符号差为 $(p, q)$ 的二阶协变对称张量场. 这意味着在每一点 $p\in\mathbf{M}$ 可以选取\newterm{tangent-space}{切空间}{tangent space} $T\mathbf{M}_p$ 中由 $p+q$ 个向量组成的基 $\{e_{(\alpha)}\}$, 使得

\begin{equation}
\mathbf{g}(e_{(\alpha)}, e_{(\beta)}) = \eta_{\alpha\beta}\label{eq:cA-401}
\end{equation}

对一切 $\alpha, \beta=0, 1, \ldots, n$ 成立, 其中 $\eta$ 是对角矩阵, 其前 $p$ 个对角元为 $-1$, 后 $q$ 个对角元为 $+1$. 若 $X$ 是 $p$ 处的任意向量, 按基 $\{e_{(\alpha)}\}$ 表为 $X=X^{\alpha}e_{(\alpha)}$, 则我们有

\begin{equation}
\mathbf{g}(X, X) = -(X^{1})^{2} - \dots - (X^{p})^{2} + (X^{p+1})^{2} + \dots + (X^{p+q})^{2}\label{eq:cA-402}
\end{equation}

当 $p=0$、$q=n$ 时, 对应于 $n$ 维 Riemann 流形. 我们感兴趣的另一个情形是 $p=1$、$q=n$, 它对应于 $n+1$ 维 Lorentz 流形. Riemann 流形的主要例子是 Euclidean 空间 $R^{n}$. 任何其他 Riemann 流形在局部上都与 $R^{n}$ 相似. 类似地, Lorentz 流形的主要例子是 Minkowski 时空, 即狭义相对论的时空. 它在 Lorentz 几何中的地位与 Euclidean 空间在 Riemann 几何中的地位相同. 此时流形 $M$ 与 $R^{n+1}$ 微分同胚, 并且存在整体定义的坐标系 $x^{\alpha}$, 相对于该坐标系度量取对角形式 $-1, 1, \ldots, 1$. 所有这些坐标系都通过 Lorentz 变换相联系, 并称为\newterm{inertial-coordinates}{惯性坐标}{inertial coordinates}. 我们将 $n+1$ 维 Minkowski 时空记作 $(\mathbb{R}^{n+1}, \mathbf{m})$.

相对于给定的坐标系 $x^{\mu}$, 伪 Riemann 度量的分量取如下形式

\[
g_{\mu\nu} = \mathbf{g}(\partial_{\mu}, \partial_{\nu})
\]

其中 $\partial_{\mu}=\frac{\partial}{\partial x^{\mu}}$ 是相应的坐标向量场. 我们用 $g^{\mu\nu}$ 表示逆度量 $g^{-1}$ 在同一坐标系 $x$ 下的分量, 用 $|g|$ 表示矩阵 $g_{\mu\nu}$ 的行列式. $\mathbf{M}$ 的体积元 $dv_{\mathbf{M}}$ 在局部坐标下表示为 $\sqrt{|g|}dx=\sqrt{|g|}dx^{1}\ldots dx^{n}$. 因此, 支集在坐标卡 $U\subset M$ 中的函数 $f$ 的积分 $\int_{M}fdv_{M}$ 定义为 $\int_{U}f(x)\sqrt{|g(x)|}dx$. 任意函数 $f$ 在 $M$ 上的积分通过作从属于坐标卡覆盖的单位分解来定义. 容易验证该定义与局部坐标系的选取无关.

由 \eqref{eq:cA-402} 可见, Lorentz 度量把每一点 $p$ 处切空间 $TM_{p}$ 中的向量分为\newterm{timelike-vector}{类时}{timelike}、\newterm{null-vector}{类光}{null}或\newterm{spacelike-vector}{类空}{spacelike}三类, 分别取决于二次型

\begin{equation}
(X, X) = g_{\mu\nu} X^{\mu} X^{\nu}\label{eq:cA-403}
\end{equation}

为负、为零或为正. 类光向量的集合 $N_{p}$ 构成一个双锥, 称为点 $p$ 的零锥. 类时向量的集合 $I_{p}$ 构成该锥的内部. $I_{p} \cup N_{p}$ 中的向量称为\newterm{causal-vector}{因果的}{causal}. 类空向量的集合 $S_{p}$ 是 $I_{p} \cup N_{p}$ 的补集.

满足 \eqref{eq:cA-401} 的标架 $e_{(\alpha)}$ 称为标准正交标架. 在 Lorentz 流形的情形下, 除了标准正交标架之外, 考虑零标架也是有意义的. 零标架是这样的向量场\footnote{只要不会引起混淆, 为简化记号我们写 $e_{\alpha}$ 而非 $e_{(\alpha)}$.}集合 $e_{\alpha}$, 它由两个类光向量 $e_{n+1}, e_{n}$ 与标准正交的类空向量 $(e_{a})_{a=1, \ldots, n-1}$ 组成, 满足

\[
\begin{array}{l}
\mathbf{g}(e_{n}, e_{n}) = \mathbf{g}(e_{n+1}, e_{n+1}) = 0, \quad \mathbf{g}(e_{n}, e_{n+1}) = -2 \\
\mathbf{g}(e_{n}, e_{a}) = \mathbf{g}(e_{n+1}, e_{a}) = 0, \quad \mathbf{g}(e_{a}, e_{b}) = \delta_{ab}
\end{array}
\]

一次形式 $A = A_{\alpha}dx^{\alpha}$ 是 $M$ 的余切丛的截面. 我们用 $A(X)$ 表示 $A$ 与向量场 $X$ 之间的自然配对. 我们可以通过 $A^{\alpha} = g^{\alpha\beta}A_{\beta}$ 提升 $A$ 的指标. $A' = A^{\alpha}\partial_{\alpha}$ 定义 $M$ 上的向量场, 且有 $A(X) = \mathbf{g}(A', X)$. k 阶协变张量 $A$ 是 TM 上的 $k$ 重线性形式.

\par\noindent\textbf{记号:}我们将使用以下记号约定:用粗体字母表示重要的张量, 如度量 $g$ 与 Riemann 曲率张量 R. 它们相对于任意标架的分量也用粗体字母表示. 因此, 给定标架 $\{e_{(\alpha)}\}$, 我们写 $\mathbf{g}_{\alpha\beta} = \mathbf{g}(e_{\alpha}, e_{\beta})$、$\mathbf{R}_{\alpha\beta\gamma\delta} = \mathbf{R}(e_{\alpha}, e_{\beta}, e_{\gamma}, e_{\delta})$, 并且对任意张量 T,

\[
T_{\alpha\beta\gamma\delta\dots} \equiv T(e_{\alpha}, e_{\beta}, e_{\gamma}, e_{\delta}, \dots)
\]

在 Riemann 流形的情形下, 我们常用拉丁字母 $i, j, k, l, \ldots$ 表示坐标 $x^{1}, x^{2}, \ldots, x^{n}$ 或张量的指标. 对 Lorentz 流形, 我们用希腊字母 $\alpha, \beta, \gamma, \ldots$ 表示指标 $0, 1, \ldots, n$. 对一般的符号差为 $s$ 的伪 Riemann 度量, 我们也使用希腊指标.

下面我们将回顾以下主题:

\begin{enumerate}
\renewcommand{\labelenumi}{\arabic{enumi}.)}
\item 向量场的 Lie 括号. Frobenius 定理
\item 张量场的 Lie 导数
\item 多重线性形式与外微分
\item 联络与协变导数
\item 伪 Riemann 度量. Riemann 与 Lorentz 几何.
\item 与伪 Riemann 度量相伴的 Levi-Civita 联络.
\item 平行移动、测地线、指数映射、完备性
\item 伪 Riemann 流形的曲率张量. 对称性. 第一与第二 Bianchi 恒等式.
\item 等距与共形等距. Killing 与共形 Killing 向量场.
\end{enumerate}


\section{各种微分概念}

这里我们回顾 Riemann 或 Lorentz 流形上微分几何的三个基本算子:外微分、\newterm{lie-derivative}{Lie 导数}{Lie derivative}, 以及 Levi-Civita 联络及其相应的\newterm{covariant-derivative}{协变导数}{covariant derivative}.

\mtsubsec{外微分}

给定标量函数 $f$, 其微分 $df$ 是由下式定义的一次形式

\[
df(X) = X(f)
\]

对任意向量场 X. 这个定义可以按如下方式推广到 $M$ 上的所有微分形式, 即 $\Omega^{k}(\mathbf{M})$:

\begin{enumerate}
\renewcommand{\labelenumi}{\roman{enumi})}
\item $d$ 是从 $\mathbf{M}$ 上所有 $k$ 形式的空间到 $k+1$ 形式的空间的线性算子. 因此对一切 $k$ 形式 $A$、$B$ 与实数 $\lambda, \mu$,
\end{enumerate}

\[
d(\lambda A + \mu B) = \lambda dA + \mu dB
\]

\begin{enumerate}
\renewcommand{\labelenumi}{\roman{enumi})}
\setcounter{enumi}{1}
\item 对任意 $k$ 形式 $A$ 与任意形式 B,
\end{enumerate}

\[
d(A \wedge B) = dA \wedge B + (-1)^{k} A \wedge dB
\]

\begin{enumerate}
\renewcommand{\labelenumi}{\roman{enumi})}
\setcounter{enumi}{2}
\item 对任意形式 A,
\end{enumerate}

\[
d^{2}A = 0.
\]

我们回忆, 若 $\Phi$ 是从 $M$ 到另一流形 $M'$ 的光滑映射, 则

\[
d(\Phi^{*}A) = \Phi^{*}(dA)
\]

最后, 若 $A$ 是一次形式而 $X$、$Y$ 是任意向量场, 我们还有方程

\[
dA(X, Y) = \frac{1}{2}\left(X(A(Y)) - Y(A(X)) - A([X, Y])\right)
\]

其中 $[X, Y]$ 是换位子 $X(Y)-Y(X)$. 这可以很容易地推广到任意的 $k$ 形式, 参见 \emph{\cite{sp}} 第 I 卷第 7 章定理 13.

\mtsubsec{Lie 导数}

考虑任意向量场 X. 在局部坐标 $x^{\mu}$ 下, X 的流由微分方程组

\[
\frac{dx^{\mu}}{dt} = X^{\mu}(x^{1}(t), \dots, x^{p+q}(t)).
\]

给出. 相应的曲线 $x^{\mu}(t)$ 称为 $X$ 的积分曲线. 对每一点 $p \in \mathbf{M}$, 存在一个开邻域 $\mathcal{U}$、一个小量 $\epsilon > 0$ 和一族微分同胚 $\Phi_t: \mathcal{U} \to \mathbf{M}$, $|t| \leq \epsilon$, 它把 $\mathcal{U}$ 中的每一点沿 $X$ 的积分曲线送到参数距离为 $t$ 处. 我们利用这些微分同胚来构造:对 $p$ 处的任意张量 T, $\Phi_t(p)$ 处的张量族 $(\Phi_t)_*T$.

张量场 $T$ 关于 $X$ 的 Lie 导数 $L_{X}T$ 为

\[
\mathcal{L}_{X}T|_{p} \equiv \lim_{t \to 0} \frac{1}{t}\left(T|_{p} - (\Phi_{t})_{*}T|_{p}\right).
\]

它具有以下性质:

\begin{enumerate}
\renewcommand{\labelenumi}{\roman{enumi})}
\item $L_{X}$ 线性地把 $(p, q)$ 型张量场映为同型的张量场.
\item $\mathcal{L}_X$ 与缩并可交换.
\item 对任意张量场 $S$、T,
\end{enumerate}

\[
\mathcal{L}_{X}(S \otimes T) = \mathcal{L}_{X}S \otimes T + S \otimes \mathcal{L}_{X}T.
\]

若 $X$ 是向量场, 容易验证

\[
\mathcal{L}_{X}Y = [X, Y]
\]

方法是写出 $(\mathcal{L}_XY)^i = -\frac{d}{dt}((\Phi_t)_*Y)^i\Big|_{t = 0}$, 并把 $((\Phi_t)_*Y)^i\Big|_p = \frac{\partial x^i(\Phi_t(q))}{\partial x^j(q)} Y^j\Big|_q$ 代入, 其中 $q = \Phi_{-t}(p)$. (详见 Hawking 与 Ellis 第 2.4 节.)

若 $A$ 是 $k$ 形式, 作为外微分与回拉 $\Phi^{*}$ 的交换公式的推论, 我们有

\[
d(\mathcal{L}_{X}A) = \mathcal{L}_{X}(dA).
\]

对给定的 $k$ 阶协变张量场 T, 我们有

\[
\mathcal{L}_{X}T(Y_{1}, \dots, Y_{k}) = XT(Y_{1}, \dots, Y_{k}) - \sum_{i=1}^{k} T(Y_{1}, \dots, \mathcal{L}_{X}Y_{i}, \dots, Y_{k})
\]

我们注意到两个坐标向量场的 Lie 括号为零,

\[
\left[\frac{\partial}{\partial x^{\mu}}, \frac{\partial}{\partial x^{\nu}}\right] = 0.
\]

其逆命题也成立, 即(参见 \emph{\cite{sp}} 第 I 卷第 5 章):

\begin{proposition}\label{prop:cA-2.3}
若 $X_{(0)}, \ldots, X_{(k)}$ 是点 $p$ 的邻域中的线性无关向量场, 且其中任意两个的 Lie 括号均为零, 则存在 $p$ 附近的坐标系 $x^{\mu}$, 使得对每个 $\rho = 0, \ldots, k$ 有 $X_{(\rho)} = \frac{\partial}{\partial x^{\rho}}$.
\end{proposition}

\begin{proof}
当 k = 2 时, 该命题等价于如下事实:若两个线性无关的向量场 $X$、$Y$ 可交换, 即 $[X, Y] = 0$, 则它们的流在如下意义下也可交换. 从任意点 $p$ 出发, 若先沿从 $p$ 出发的 $X$ 的积分曲线行进参数距离 t, 再沿 $Y$ 的积分曲线行进参数距离 s, 则到达的点与先沿 $Y$ 的积分曲线行进距离 $s$、再沿 $X$ 的积分曲线行进距离 $t$ 所到达的点相同. 为了证明该断言, 我们首先考虑相对于 $X = \partial_{t}$ 的坐标 t, x. Y 在同一坐标系下的系数与 $t$ 无关. 然后作另一个坐标变换 $x = x(y)$, 使得 $Y = A(y)\partial_{t} + \partial_{y^{1}}$. 最后可以再作形如 $s = a(y) + t$ 的坐标变换, 其中 $a$ 是 $A$ 关于 $y^{1}$ 的原函数. 在新坐标 s, y 下, 必须有 $X = \partial_{s}$, $Y = \partial_{y^{1}}$, 正如所愿. 一般情形可以同样证明.
\end{proof}

上述命题是证明 Frobenius 定理的主要步骤. 为了陈述该定理, 我们回忆 $M$ 中 $k$ 维\newterm{k-distribution}{分布}{k-distribution}的定义. 所谓 $k$ 维分布, 是指在 $U$ 的定义域中每一点处光滑地指定一个 $k$ 维平面 $\pi_{p}$. 若对 $U$ 上的任意向量场 $X$、Y, 只要对任意 $p \in U$ 有 $X|_{p}, Y|_{p} \in \pi_{p}$, 就有 $[X, Y]|_{p} \in \pi_{p}$, 则该分布称为\newterm{involutive}{对合的}{involutive}. 对\newterm{integrable}{可积的}{integrable}分布\footnote{回忆:$U$ 上的分布 $\pi$ 称为可积的, 如果通过每一点 $p \in U$ 恰有一条维数为 $k$ 的子流形 N, 使得 $\pi_p = TN_p$.}显然如此. 事实上, 若对一切 $p \in N$ 有 $X|_{p}, Y|_{p} \in TN_{p}$, 则 $X$、$Y$ 切于 N, 它们的换位子 $[X, Y]$ 亦然. Frobenius 定理断言其逆命题也成立\footnote{证明参见 Spivak 的《[Sp]》第 I 卷第 6 章.}, 即处于对合状态也是分布可积的充分条件:

\begin{theorem}[Frobenius 定理]\label{thm:cA-2.4}
分布 $(\pi_p)_{p\in\mathcal{U}}$ 可积的充分必要条件是它是对合的.
\end{theorem}

\begin{proof}
若 $\pi$ 是对合的, 我们可以找到线性无关的向量场 $X_{1}, \ldots, X_{k}$, 它们在每一点张成 $\pi$, 并且其中任意两个的 Lie 括号都是 $X_{1}, \ldots, X_{k}$ 的线性组合. 然后可以重新定义 $X_{1}^{\prime}$ 为其他向量场的线性组合, 使得它与 $X_{2}, \ldots, X_{k}$ 可交换. 接着可以归纳地重新定义其余向量场, 使它们两两可交换. 证明于是由前一命题得出.
\end{proof}

\mtsubsec{联络与协变导数}

\newterm{connection}{联络}{connection} D 是这样一种法则:它给每个向量场 $X$ 指派一个微分算子 $D_{X}$. 该算子把向量场 $Y$ 映为向量场 $D_{X}Y$, 且对 $\alpha, \beta \in R$ 与 $M$ 上的标量函数 $f$、g, 满足

\begin{equation}
\begin{array}{l}
a) \quad \mathbf{D}_{fX + gY}Z = f\mathbf{D}_{X}Z + g\mathbf{D}_{Y}Z \\
b) \quad \mathbf{D}_{X}(\alpha Y + \beta Z) = \alpha\mathbf{D}_{X}Y + \beta\mathbf{D}_{X}Z \\
c) \quad \mathbf{D}_{X}fY = X(f)Y + f\mathbf{D}_{X}Y
\end{array}\label{eq:cA-404}
\end{equation}

因此, 在点 $p$ 处,

\begin{equation}
\mathbf{D}Y \equiv Y^{\alpha}_{;\beta}\, \theta^{(\beta)} \otimes e_{(\alpha)}\label{eq:cA-405}
\end{equation}

其中 $\theta^{(\beta)}$ 是关于标准正交标架 $e_{(\beta)}$ 的对偶基的一次形式. 注意 $Y_{;\beta}^{\alpha} = \theta^{(\alpha)}(\mathbf{D}_{e_{(\beta)}}Y)$. 另一方面, 由 c) 得

\[
\mathbf{D}(fY) = df \otimes Y + f\mathbf{D}Y
\]

所以

\[
\mathbf{D}Y = \mathbf{D}(Y^{\alpha}e_{(\alpha)}) = dY^{\alpha} \otimes e_{(\alpha)} + Y^{\alpha}\mathbf{D}e_{(\alpha)}
\]

最后, 利用 $df(\cdot) = e_{(\alpha)}(f)\theta^{(\alpha)}(\cdot)$, 得

\begin{equation}
\mathbf{D}Y = \left(e_{(\beta)}(Y^{\alpha}) + Y^{\gamma}\theta^{(\alpha)}(\mathbf{D}_{e_{(\beta)}}e_{(\gamma)})\right)\theta^{(\beta)} \otimes e_{(\alpha)}\label{eq:cA-406}
\end{equation}

因此

\[
Y^{\alpha}_{;\beta} = e_{(\beta)}(Y^{\alpha}) + \Gamma^{\alpha}_{\beta\gamma}Y^{\gamma}
\]

从而联络由其联络系数

\begin{equation}
\Gamma^{\alpha}_{\beta\gamma} = \theta^{(\alpha)}(\mathbf{D}_{e_{(\beta)}}e_{(\gamma)})\label{eq:cA-407}
\end{equation}

决定, 这些联络系数在坐标基下就是通常的Christoffel 符号, 具有表达式

\[
\Gamma^{\mu}_{\rho\nu} = dx^{\mu}\left(\mathbf{D}_{\frac{\partial}{\partial x^{\rho}}}\frac{\partial}{\partial x^{\nu}}\right)
\]

最后

\begin{equation}
\mathbf{D}_{X}Y = \left(X(Y^{\alpha}) + \Gamma^{\alpha}_{\beta\gamma}X^{\beta}Y^{\gamma}\right)e_{(\alpha)}\label{eq:cA-408}
\end{equation}

在坐标标架的特殊情形下, 我们有

\[
\mathbf{D}_{X}Y = \left(X^{\mu}\frac{\partial Y^{\nu}}{\partial x^{\mu}} + \Gamma^{\nu}_{\rho\sigma}X^{\rho}Y^{\sigma}\right)\frac{\partial}{\partial x^{\nu}}
\]

若 $\mathbf{D}\mathbf{g} = 0$, 则联络称为\newterm{levi-civita-connection}{Levi-Civita 联络}{Levi-Civita connection}. 也就是说, 对任意三个向量场 X, Y, Z,

\begin{equation}
Z(\mathbf{g}(X, Y)) = \mathbf{g}(\mathbf{D}_{Z}X, Y) + \mathbf{g}(X, \mathbf{D}_{Z}Y)\label{eq:cA-409}
\end{equation}

微分几何中一个非常简单而基本的结果断言:对任意给定的(伪 Riemann)度量, 存在唯一与其相伴的仿射联络.

\begin{proposition}\label{prop:cA-2.6}
在 $M$ 上存在唯一的联络, 称为 Levi-Civita 联络, 它满足 Dg = 0. 该联络是无挠的, 即

\[
\mathbf{D}_{X}Y - \mathbf{D}_{Y}X = [X, Y].
\]

此外, 相对于坐标系 $x^{\mu}$, 该联络的 Christoffel 符号由标准公式

\[
\Gamma^{\mu}_{\rho\nu} = \frac{1}{2}g^{\mu\tau}\left(\partial_{\rho}g_{\nu\tau} + \partial_{\nu}g_{\tau\rho} - \partial_{\tau}g_{\nu\rho}\right).
\]

给出.
\end{proposition}

\par\noindent\textbf{习题:}自己证明上述命题, 不要看书.

到目前为止, 我们只定义了向量场的协变导数. 我们可以容易地把定义推广到一次形式 $A = A_{\alpha}dx^{\alpha}$:要求对一切向量场 $X$、$Y$ 有

\[
X(A(Y)) = \mathbf{D}_{X}A(Y) + A(D_{X}Y),
\]

给定 $k$ 阶协变张量场 T, 我们按如下法则定义其协变导数 $D_{X}T$:

\[
\mathbf{D}_{X}T(Y_{1}, \dots, Y_{k}) = XT(Y_{1}, \dots, Y_{k}) - \sum_{i=1}^{k} T(Y_{1}, \dots, \mathbf{D}_{X}Y_{i}, \dots, Y_{k})
\]

我们可以把 DT 看作秩为 $k+1$ 的协变张量, 定义为

\[
\mathbf{D}T(X, Y_{1}, \dots, Y_{k}) = \mathbf{D}_{X}T(Y_{1}, \dots, Y_{k}).
\]

给定标架 $e_{\alpha}$, 我们用 $T_{\alpha_{1}\ldots, \alpha_{k};\beta} = \mathbf{D}T(e_{\beta}, e_{\alpha_{1}}, \ldots, e_{\alpha_{k}})$ 表示 DT 相对于该标架的分量. 通过重复的协变微分, 我们可以定义 $D^{2}T, \ldots D^{m}T$. 相对于标架 $e_{\alpha}$, 我们写

\[
\mathbf{D}_{\beta_{1}}\dots\mathbf{D}_{\beta_{m}}T_{\alpha_{1}\dots\alpha_{k}} = T_{\alpha_{1}\dots\alpha_{k};\beta_{1}\dots\beta_{m}} = \mathbf{D}^{m}T(e_{\beta_{1}}\dots, e_{\beta_{m}}, e_{\alpha_{1}}, \dots, e_{\alpha_{k}}).
\]

Levi-Civita 联络无挠这一事实使我们能把协变微分与 Lie 导数联系起来. 于是, 若 $T$ 是 $k$ 阶协变张量, 在坐标基下我们有

\[
(\mathcal{L}_{X}T)_{\sigma_{1}\ldots\sigma_{k}} = X^{\mu}T_{\sigma_{1}\ldots\sigma_{k};\mu} + X^{\mu}_{;\sigma_{1}}T_{\mu\sigma_{2}\ldots\sigma_{k}} + \ldots + X^{\mu}_{;\sigma_{k}}T_{\sigma_{1}\ldots\sigma_{k-1}\mu}.
\]

协变导数还与外导数按如下简单公式相联系. 若 $A$ 是 $k$ 形式, 我们有\footnote{$[\sigma_1\ldots\sigma_k;\mu]$ 表示对所有指标取反对称化(即张量在所有指标置换上的交错和乘以 $1/k!$), 而 "$, \mu$" 表示对 $x^{\mu}$ 的普通偏导数.}$A_{[\sigma_{1}\ldots\sigma_{k};\mu]} = A_{[\sigma_{1}\ldots\sigma_{k}, \mu]}$, 并且

\[
dA = \sum A_{\sigma_{1}\dots\sigma_{k};\mu}\, dx^{\mu} \wedge dx^{\sigma_{1}} \wedge dx^{\sigma_{2}} \wedge \dots \wedge dx^{\sigma_{k}}.
\]

给定光滑曲线 $x : [0, 1] \to M$, 以 $t$ 为参数, 令 $T = \left(\frac{\partial}{\partial t}\right)_{\mathbf{x}}$ 为沿曲线的相应切向量场. 定义在曲线上的向量场 $X$ 称为沿曲线\newterm{parallel-transport}{平行移动}{parallel transport}的, 如果 $D_{T}X = 0$. 若曲线相对于坐标系有参数方程 $x^{\nu} = x^{\nu}(t)$, 则 $T^{\mu} = \frac{dx^{\mu}}{dt}$, 并且分量 $X^{\mu} = X^{\mu}(\mathbf{x}(t))$ 满足常微分方程组

\[
\frac{\mathbf{D}}{dt}X^{\mu} \equiv \frac{dX^{\mu}}{dt} + \Gamma^{\mu}_{\rho\sigma}(\mathbf{x}(t))\frac{dx^{\rho}}{dt}X^{\sigma} = 0.
\]

若在曲线的每一点处 $D_{T}T$ 切于曲线, 即 $D_{T}T = \lambda T$, 则称该曲线为\newterm{geodesic}{测地线}{geodesic}. 此时可以重新参数化曲线, 使得相对于新参数 s, 切向量 $S = \left(\frac{\partial}{\partial s}\right)_{\mathbf{x}}$ 满足 $D_{S}S = 0$. 这样的参数称为"仿射参数". 仿射参数在变换 $s = as' + b$(a、$b$ 为常数)意义下确定. 相对于仿射参数 $s$ 与任意坐标 $x^{\mu}$, 测地线满足方程

\[
\frac{d^{2}x^{\mu}}{ds^{2}} + \Gamma^{\mu}_{\rho\sigma}\frac{dx^{\rho}}{ds}\frac{dx^{\sigma}}{ds} = 0.
\]

以仿射参数参数化的测地曲线简称为测地线. 在 Lorentz 几何中, 类时测地线对应于在由联络系数表示的引力场中自由下落的粒子的\newterm{world-line}{世界线}{world line}. 此时仿射参数 $s$ 称为粒子的\newterm{proper-time}{固有时}{proper time}.

给定点 $p \in M$ 与切空间 $T_{p}M$ 中的向量 X, 令 $\mathbf{x}(t)$ 为从 $p$ 出发、具有"速度" X 的唯一测地线. 我们定义\newterm{exponential-map}{指数映射}{exponential map}:

\[
\exp_{p}: T_{p}\mathbf{M} \to \mathbf{M}.
\]

这个映射不一定对所有 $X \in T_{p}M$ 都有定义. 常微分方程组的存在唯一性定理表明指数映射在 $T_{p}M$ 中原点的一个邻域上有定义. 若对每一点 p, 指数映射对整个 $T_{p}M$ 都有定义, 则称流形 $M$ 是\newterm{geodesically-complete}{测地完备的}{geodesically complete}. 一般而言, 若联络是 $C^{r}$ 联络\footnote{所谓 $C^{r}$ 联络是指:若 $Y$ 是 $C^{r+1}$ 向量场, 则 DY 是 $C^{r}$ 向量场.}, 则存在 $T_{p}M$ 中原点的开邻域 $U_{0}$ 与 $M$ 中点 $p$ 的开邻域 $V_{p}$, 使得映射 $\exp_{p}$ 是 $U_{0}$ 到 $V_{p}$ 上的 $C^{r}$ 微分同胚. 邻域 $V_{p}$ 称为 $p$ 的\newterm{normal-neighbourhood}{法邻域}{normal neighbourhood}.


\section{Riemann 曲率张量、Ricci 张量与 Bianchi 恒等式}

在平坦时空中, 若沿任意闭曲线平行移动一个向量, 则得到出发时的那个向量. 一般而言这不再成立, 因为向量场的二阶协变导数不可交换. 这种不可交换性由\newterm{curvature-tensor}{曲率张量}{curvature tensor}度量:

\begin{equation}
\mathbf{R}(X, Y)Z = \mathbf{D}_{X}(\mathbf{D}_{Y}Z) - \mathbf{D}_{Y}(\mathbf{D}_{X}Z) - \mathbf{D}_{[X, Y]}Z\label{eq:cA-410}
\end{equation}

或者写成相对于任意标架的分量形式:

\begin{equation}
\mathbf{R}^{\alpha}_{\beta\gamma\delta} = \theta^{(\alpha)}\big((\mathbf{D}_{\gamma}\mathbf{D}_{\delta} - \mathbf{D}_{\delta}\mathbf{D}_{\gamma})e_{(\beta)}\big)\label{eq:cA-411}
\end{equation}

相对于坐标系 $x^{\mu}$, 用 $g_{\mu\nu}$ 分量写出, Riemann 分量的表达式为

\begin{equation}
R^{\mu}_{\nu\rho\sigma} = \frac{\partial\Gamma^{\mu}_{\sigma\nu}}{\partial x^{\rho}} - \frac{\partial\Gamma^{\mu}_{\rho\nu}}{\partial x^{\sigma}} + \Gamma^{\mu}_{\rho\tau}\Gamma^{\tau}_{\sigma\nu} - \Gamma^{\mu}_{\sigma\tau}\Gamma^{\tau}_{\rho\nu}\label{eq:cA-412}
\end{equation}

曲率张量的基本性质(由 Riemann 首先证明)断言:若 $\mathbf{R}$ 在点 $p$ 的邻域中恒为零, 则可以找到一族局部坐标, 使得在 $p$ 的邻域中 $g_{\mu\nu} = \eta_{\mu\nu}$\footnote{关于此事实的详细讨论与证明, 参见 Spivak 的《\cite{sp}》第 II 卷.}.

曲率张量关于度量 $g$ 的迹是一个对称张量, 称为\newterm{ricci-tensor}{Ricci 张量}{Ricci tensor},

\[
\mathbf{R}_{\alpha\beta} = \mathbf{g}^{\gamma\delta}\mathbf{R}_{\alpha\gamma\beta\delta}
\]

标量曲率是 Ricci 张量的迹

\[
\mathbf{R} = \mathbf{g}^{\alpha\beta}\mathbf{R}_{\alpha\beta}.
\]

任意时空 $(\mathbf{M}, \mathbf{g})$ 的 Riemann 曲率张量具有以下对称性质,

\begin{equation}
\begin{aligned}
& \mathbf{R}_{\alpha\beta\gamma\delta} = -\mathbf{R}_{\beta\alpha\gamma\delta} = -\mathbf{R}_{\alpha\beta\delta\gamma} = \mathbf{R}_{\gamma\delta\alpha\beta} \\
& \mathbf{R}_{\alpha\beta\gamma\delta} + \mathbf{R}_{\alpha\gamma\delta\beta} + \mathbf{R}_{\alpha\delta\beta\gamma} = 0
\end{aligned}\label{eq:cA-413}
\end{equation}

\eqref{eq:cA-413} 中的第二个恒等式称为第一 Bianchi 恒等式.

它还满足第二 Bianchi 恒等式, 我们在这里称之为 Bianchi 方程, 在一般标架下具有如下形式:

\begin{equation}
\mathbf{D}_{[\epsilon}\mathbf{R}_{\gamma\delta]\alpha\beta} = 0\label{eq:cA-414}
\end{equation}

曲率张量的无迹部分 $\mathbf{C}$ 称为\newterm{weyl-tensor}{Weyl 张量}{Weyl tensor}, 在任意标架下具有如下表达式:

\begin{equation}
\begin{aligned}
\mathbf{C}_{\alpha\beta\gamma\delta} = &\ \mathbf{R}_{\alpha\beta\gamma\delta} - \frac{1}{n-1}\left(\mathbf{g}_{\alpha\gamma}\mathbf{R}_{\beta\delta} + \mathbf{g}_{\beta\delta}\mathbf{R}_{\alpha\gamma} - \mathbf{g}_{\beta\gamma}\mathbf{R}_{\alpha\delta} - \mathbf{g}_{\alpha\delta}\mathbf{R}_{\beta\gamma}\right) \\
& + \frac{1}{n(n-1)}(\mathbf{g}_{\alpha\gamma}\mathbf{g}_{\beta\delta} - \mathbf{g}_{\alpha\delta}\mathbf{g}_{\beta\gamma})\mathbf{R}
\end{aligned}\label{eq:cA-415}
\end{equation}

注意 $C$ 满足 Riemann 张量的所有对称性质:

\begin{equation}
\begin{aligned}
& \mathbf{C}_{\alpha\beta\gamma\delta} = -\mathbf{C}_{\beta\alpha\gamma\delta} = -\mathbf{C}_{\alpha\beta\delta\gamma} = \mathbf{C}_{\gamma\delta\alpha\beta} \\
& \mathbf{C}_{\alpha\beta\gamma\delta} + \mathbf{C}_{\alpha\gamma\delta\beta} + \mathbf{C}_{\alpha\delta\beta\gamma} = 0
\end{aligned}\label{eq:cA-416}
\end{equation}

此外还有 $\mathbf{g}^{\alpha\gamma}\mathbf{C}_{\alpha\beta\gamma\delta} = 0$.

我们说两个度量 $g$ 与 $\hat{g}$ 是\newterm{conformal}{共形的}{conformal}的, 如果存在某个非零可微函数 $\lambda$ 使得 $\hat{g} = \lambda^{2}g$. 于是下面的定理成立(参见 \emph{Hawking--Ellis}, 第 2 章第 2.6 节):

\begin{theorem}\label{thm:cA-3.1}
设 $\hat{\mathbf{g}} = \lambda^2\mathbf{g}$, $\hat{\mathbf{C}}$ 是相对于 $\hat{\mathbf{g}}$ 的 Weyl 张量, $\mathbf{C}$ 是相对于 $\mathbf{g}$ 的 Weyl 张量. 则

\[
\hat{\mathbf{C}}^{\alpha}_{\beta\gamma\delta} = \mathbf{C}^{\alpha}_{\beta\gamma\delta}.
\]
\end{theorem}

因此 $\mathbf{C}$ 是\newterm{conformally-invariant}{共形不变的}{conformally invariant}.

\mtsubsec{等距与共形等距、Killing 与共形 Killing 向量场}

\begin{definition*}
微分同胚 $\Phi : \mathcal{U} \subset \mathbf{M} \to \mathbf{M}$ 称为\newterm{conformal-isometry}{共形等距}{conformal isometry}, 如果在每一点 $p$ 有 $\Phi_{*}\mathbf{g} = \Lambda^{2}\mathbf{g}$, 即

\[
(\Phi^{*}\mathbf{g})(X, Y)|_{p} = \mathbf{g}(\Phi_{*}X, \Phi_{*}Y)|_{\Phi(p)} = \Lambda^{2}\mathbf{g}(X, Y)|_{p}
\]

其中 $\Lambda\neq 0$. 若 $\Lambda = 1$, 则称 $\Phi$ 为 $M$ 的等距.
\end{definition*}

\begin{definition*}
生成单参数等距(相应地, 共形等距)群的向量场 $K$ 称为\newterm{killing-vector-field}{Killing 向量场}{Killing vector field}(相应地, \newterm{conformal-killing-vector-field}{共形 Killing 向量场}{conformal Killing vector field}).
\end{definition*}

设 $K$ 是这样的向量场, $\Phi_{t}$ 是相应的单参数群. 由于 $(\Phi_{t})_{*}$ 是共形等距, 我们推得 $L_{K}g$ 必与度量 $g$ 成比例. 此外, 若 $K$ 是 Killing 向量场, 则 $L_{K}g = 0$.

\begin{definition*}
给定任意向量场 X, 我们用 $^{(X)}\pi$ 表示 $X$ 的变形张量, 由如下公式定义:

\[
{}^{(X)}\pi_{\alpha\beta} = (\mathcal{L}_{X}g)_{\alpha\beta} = \mathbf{D}_{\alpha}X_{\beta} + \mathbf{D}_{\beta}X_{\alpha}.
\]
\end{definition*}

张量 $(X)\pi$ 在精确的意义上度量了由 $X$ 生成的微分同胚与等距或共形等距相差多少. 下面的命题成立(参见 \emph{Hawking--Ellis}, 第 2 章第 2.6 节):

\begin{proposition}\label{prop:cA-3.3}
向量场 $X$ 是 Killing 的当且仅当 $(X)\pi = 0$. 它是共形 Killing 的当且仅当 $(X)\pi$ 与 $g$ 成比例.
\end{proposition}

\begin{remark*}
可以选择局部坐标使得 $X = \frac{\partial}{\partial x^{\mu}}$. 于是立即推得, 相对于这些坐标, 度量 $g$ 与分量 $x^{\mu}$ 无关.
\end{remark*}

\begin{proposition}\label{prop:cA-3.4}
在任意维数为 $n = p + q$ 的伪 Riemann 时空 $M$ 上, 线性无关的 Killing 向量场不可能超过 $\frac{1}{2}(p + q)(p + q + 1)$ 个.
\end{proposition}

\begin{proof}
\hyperref[prop:cA-3.4]{\propref{prop:cA-3.4}} 是以下关系的简单推论, 该关系对任意向量场 $X$ 成立, 可由直接计算并利用 $R$ 的对称性得到:

\begin{equation}
\mathbf{D}_{\beta}\mathbf{D}_{\alpha}X_{\lambda} = \mathbf{R}_{\lambda\alpha\beta\delta}X^{\delta} + {}^{(X)}\Gamma_{\alpha\beta\lambda}\label{eq:cA-417}
\end{equation}

其中

\begin{equation}
{}^{(X)}\Gamma_{\alpha\beta\lambda} = \frac{1}{2}\left(\mathbf{D}_{\beta}\pi_{\alpha\lambda} + \mathbf{D}_{\alpha}\pi_{\beta\lambda} - \mathbf{D}_{\lambda}\pi_{\alpha\beta}\right)\label{eq:cA-418}
\end{equation}

而 $\pi \equiv {}^{(X)}\pi$ 是 $X$ 的变形张量.

若 $X$ 是 Killing 向量场, 方程 \eqref{eq:cA-417} 变为

\begin{equation}
\mathbf{D}_{\beta}(\mathbf{D}_{\alpha}X_{\lambda}) = \mathbf{R}_{\lambda\alpha\beta\delta}X^{\delta}\label{eq:cA-419}
\end{equation}

再结合常微分方程的存在唯一性定理, 这意味着任何 Killing 向量场完全由 $X$ 与 DX 在给定点的 $\frac{1}{2}(n+1)(n+2)$ 个值确定. 事实上, 设 $p$、$q$ 是由具有切向量 $T$ 的曲线 $x(t)$ 相连的两点. 令 $L_{\alpha\beta} \equiv D_{\alpha}X_{\beta}$, 注意沿 $x(t)$, X、$L$ 满足微分方程组

\[
\frac{\mathbf{D}}{dt}X = T \cdot L, \qquad \frac{\mathbf{D}}{dt}L = \mathbf{R}(\cdot, \cdot, X, T)
\]

因此 $X$、$L$ 沿曲线的值由它们在 $p$ 处的值唯一确定.
\end{proof}

具有最多 Killing 向量场个数的 $n$ 维 Riemann 流形是 Euclidean 空间 $R^{n}$. 类似地, Minkowski 时空 $R^{n+1}$ 是具有最多 Killing 向量场个数的 Lorentz 流形.

\mtsubsec{体积形式、Hodge 对偶与散度定理}

设 $(\mathbf{M}, \mathbf{g})$ 是维数为 $n$、符号差为 $s$ 的可定向伪 Riemann 流形 M. 设 $e_{(\alpha)}$ 是 $M$ 上任意正定向标架, $e^{(\alpha)}$ 是相应的对偶标架. 我们可以把它与 $n$ 形式

\[
e^{(1)} \wedge e^{(2)} \wedge \ldots \wedge e^{(n)}
\]

相联系. 若 $e_{(\alpha)}' = M_{\alpha}^{\beta}e_{(\beta)}$ 表示另一个基, 而 $e'^{(\alpha)} = (M^{-1})_{\beta}^{\alpha}e^{(\beta)}$ 是相应的一次形式对偶基, 则

\[
e'^{(1)} \wedge e'^{(2)} \wedge \dots \wedge e'^{(n)} = \det(M^{-1})e^{(1)} \wedge e^{(2)} \wedge \dots \wedge e^{(n)}
\]

另一方面, 若用 $\mathbf{g}$ 与 $\mathbf{g}'$ 表示矩阵 $\mathbf{g}(e_{(\alpha)}, e_{(\beta)})$、$\mathbf{g}(e'_{(\alpha)}, e'_{(\beta)})$ 的行列式, 则 $\mathbf{g}' = (\det M)^2\mathbf{g}$. 因此, n 形式

\begin{equation}
\varepsilon = |\mathbf{g}|^{1/2}e^{(1)} \wedge \dots \wedge e^{(n)}\label{eq:cA-420}
\end{equation}

与保持定向的基变换无关. 于是, 若 $\mathbf{M}$ 可定向, 我们可以定义一个唯一的 $n$ 形式, 称为 $\mathbf{M}$ 的\newterm{volume-form}{体积形式}{volume form}. 相对于正定向标准正交标架 $e_{(\alpha)}$, 当有两个指标相同时 $\varepsilon_{\alpha_1\dots\alpha_n} = 0$, 否则等于置换 $\alpha_{1}\alpha_{2}\ldots\alpha_{n}$ 的符号. $\varepsilon$ 的以下乘法性质在计算中非常有用,

\begin{equation}
\varepsilon^{\alpha_{1}\dots\alpha_{n}}\varepsilon_{\beta_{1}\dots\beta_{n}} = (-1)^{s}\det(\delta^{\alpha_{i}}_{\beta_{j}})_{i, j = 1, \dots, n}\label{eq:cA-421}
\end{equation}

其中 $s$ 是度量的符号差. 另外, 通过重复缩并,

\[
\begin{aligned}
\varepsilon^{\alpha_{1}\alpha_{2}\ldots\alpha_{n}}\varepsilon_{\alpha_{1}\beta_{2}\ldots\beta_{n}} &= (-1)^{s}\det(\delta^{\alpha_{i}}_{\beta_{j}})_{i, j = 2, \ldots, n} \\
\varepsilon^{\alpha_{1}\alpha_{2}\alpha_{3}\ldots\alpha_{n}}\varepsilon_{\alpha_{1}\alpha_{2}\beta_{3}\ldots\beta_{n}} &= 2(-1)^{s}\det(\delta^{\alpha_{i}}_{\beta_{j}})_{i, j = 3, \ldots, n} \\
\ldots &= \ldots \\
\varepsilon^{\alpha_{1}\ldots\alpha_{n}}\varepsilon_{\alpha_{1}\ldots\beta_{n}} &= n!(-1)^{s}
\end{aligned}
\]

注意 $\varepsilon$ 的协变导数消失. 事实上, 这可由对公式 $\varepsilon^{\alpha_{1}\cdots\alpha_{n}}\varepsilon_{\alpha_{1}\cdots\alpha_{n}} = n!(-1)^{s}$ 求导直接推出, 由此得 $\varepsilon^{\alpha_{1}\cdots\alpha_{n}}\varepsilon_{\alpha_{1}\cdots\alpha_{n};\mu} = 0$, 于是利用 \eqref{eq:cA-421} 有 $D_{\mu}\varepsilon_{\alpha_{1}\cdots\alpha_{n}} = \varepsilon_{\alpha_{1}\cdots\alpha_{n};\mu} = 0$.

Hodge 对偶是自然定义的从 $\Omega^{k}(\mathbf{M})$ 到 $\Omega^{n-k}(\mathbf{M})$ 的算子. 在固定的正定向标架中, 若用 $A_{\alpha_{0}\ldots\alpha_{k-1}}$ 表示 $A \in \Omega^{k}(\mathbf{M})$ 的系数, 则我们有

\[
\ast A_{\alpha_{1}\dots\alpha_{n-k}} = \frac{1}{k!}\varepsilon_{\alpha_{1}\dots\alpha_{n-k}}{}^{\beta_{1}\dots\beta_{k}}A_{\beta_{1}\dots\beta_{k}}
\]

容易验证, 对 $k$ 形式 $A$ 有\footnote{特别地, 若 $\mathbf{g}$ 是维数为 4 的 Lorentz 流形, 则 $*(*A) = (-1)^{k+1}A$.},

\begin{equation}
*(*A) = (-1)^{s}(-1)^{k(n-k)}A\label{eq:cA-422}
\end{equation}

我们按如下公式定义散度算子 $\delta : \Omega^{k}(\mathbf{M}) \to \Omega^{k-1}(\mathbf{M})$:

\begin{equation}
\delta A = *d*A\label{eq:cA-423}
\end{equation}

给定 $\mathbf{M}$ 上连续、紧支集的函数 f, 我们按如下公式定义其积分:

\[
\int_{\mathbf{M}}fdv_{\mathbf{g}} := \int f\varepsilon
\]

若 $f$ 的支集在(正定向)坐标卡内, 则

\[
\int_{\mathbf{M}}fdv_{\mathbf{g}} = \int f(x^{1}, \ldots, x^{n})|\mathbf{g}|^{1/2}dx^{1} \wedge \ldots \wedge dx^{n}.
\]

于是我们可以引入 $M$ 上函数的标量积,

\begin{equation}
\langle f_{1}, f_{2}\rangle = \int_{\mathbf{M}}f_{1}\overline{f_{2}}\, dv_{\mathbf{g}}\label{eq:cA-424}
\end{equation}

给定两个(紧支集的)形式 $A \in \Omega^k(\mathbf{M})$、$B \in \Omega^{n-k}(\mathbf{M})$, 我们定义

\[
\int_{\mathbf{M}}A \wedge B
\]

其中积分理解为 $M$ 上 $n$ 形式的通常积分. 另外, 给定 $A, B \in \Omega^{k}(\mathbf{M})$, 相对于固定的标架 $e_{1}, e_{2}$ 我们定义

\[
\mathbf{g}(A, B) = \frac{1}{k!}A^{\alpha_{1}\dots\alpha_{k}}B_{\alpha_{1}\dots\alpha_{k}}
\]

并且

\begin{equation}
\langle A, B\rangle := \int_{\mathbf{M}}\mathbf{g}(A, B)\label{eq:cA-425}
\end{equation}

利用这些定义, 我们容易验证, 对 $A \in \Omega^k(\mathbf{M})$、$B \in \Omega^{n-k}(\mathbf{M})$ 有

\[
\langle {}^{\ast}A, B\rangle = \int_{\mathbf{M}}A \wedge B
\]

还注意到散度算子 $\delta$ 在如下意义下是 $d$ 的对偶(对 $A \in \Omega^{k}(\mathbf{M})$、$B \in \Omega^{k-1}(\mathbf{M})$):

\[
\langle \delta A, B\rangle = \langle A, dB\rangle
\]

若 $X$ 是 $M$ 上的(紧支集)向量场, 我们定义其 Hodge 对偶 $*X$ 为与其相伴的一次形式的 Hodge 对偶, 即

\[
(\ast X)_{\alpha_{1}\ldots\alpha_{n-1}} = \varepsilon_{\alpha_{1}\ldots\alpha_{n-1}\mu}X^{\mu}
\]

容易验证, 相对于任意正定向坐标标架, 即 $*X = \frac{1}{(k-1)!}\varepsilon_{\alpha_{1}\ldots\alpha_{n-1}\mu}X^{\mu}dx^{\alpha_{1}} \wedge \ldots \wedge dx^{\alpha_{n-1}}$, 我们有

\[
d(*X) = (-1)^{n-1}(\mathbf{D}^{\alpha}X_{\alpha})dv_{\mathbf{g}}
\]

从而

\[
\delta X =: \mathbf{Div}X = \mathbf{D}^{\alpha}X_{\alpha}
\]

将这一事实与 Stokes 定理相结合, 我们推得如下结论.

\begin{proposition}\label{prop:cA-3.6}
设 $X$ 是 $M$ 上的向量场, $\ast X$ 是其 Hodge 对偶. 设 $U$ 是 $M$ 中具有光滑边界、且按 $M$ 诱导的正定向定向的区域. 则

\[
\int_{\mathcal{U}}\mathbf{Div}X\, dv_{\mathbf{g}} = (-1)^{n-1}\int_{\partial\mathcal{U}}{}^{\ast}X.
\]
\end{proposition}

我们将该命题应用于维数为 $n+1$ 的 Lorentz 流形. 给定具有诱导 Riemann 度量 $g$ 的类空超曲面 $\Sigma$, 我们用 $dv_{g}$ 表示 $g$ 相对于诱导正定向的体积元. 现在我们把\newterm{divergence-theorem}{散度定理}{divergence theorem}应用于 $M$ 中由两个类空超曲面 $\Sigma_{0}, \Sigma_{1}$ 围成的\newterm{lens-shaped-domain}{透镜形区域}{lens shaped domain} U.

\begin{corollary}\label{cor:cA-3.7}
设 $U$ 是由两个类空超曲面 $\Sigma_{0}, \Sigma_{1}$ 围成的透镜形区域. 设 $T$ 表示边界两侧的未来定向单位法向量, X 是 $M$ 上的任意光滑向量场. 则

\[
\int_{\Sigma_{1}}(X, T)dv_{g} - \int_{\Sigma_{0}}(X, T)dv_{g} = -\int_{\mathcal{U}}\mathbf{Div}X\, dv_{\mathbf{g}}.
\]
\end{corollary}

\begin{proof}
首先考虑边界的 $\Sigma_{0}$ 部分, 并在 $\Sigma_{0}$ 的每一点 $p$ 附近局部定义坐标系 $x^{0}, x^{1}, \ldots, x^{n}$, 使得

\begin{enumerate}
\renewcommand{\labelenumi}{(\arabic{enumi})}
\item 在 $\Sigma_0$ 上 $x^0 = 0$, 在 $\mathcal{U}$ 中 $x^0 < 0$.
\item $x^{1}, x^{2}, \ldots, x^{n}$ 在 $\Sigma_{0}$ 中是正定向的.
\item 在 $p$ 附近 $\Sigma_0$ 上 $\mathbf{g}_{00} = -1$、$\mathbf{g}_{0i} = 0$.
\end{enumerate}

相对于这些坐标, 我们有 $dv_{g} = \varepsilon_{12\ldots n}dx^{1} \wedge \ldots \wedge dx^{n}$, 其中 $\varepsilon_{12\ldots n} = \varepsilon_{01\ldots n}$. 因此在 $\Sigma_{0}$ 上,

\[
(\ast X)_{i_{1}\dots i_{n}} = \varepsilon_{i_{1}\dots i_{n}\mu}X^{\mu} = \varepsilon_{i_{1}\dots i_{n}0}X^{0} = (-1)^{n}X^{0}dv_{g}
\]
\end{proof}

\mtsubsec{Laplace--Beltrami 算子}

伪 Riemann 流形 $M$ 上的标量Laplace--Beltrami 算子定义为

\begin{equation}
\Delta_{\mathbf{M}}u(x) = g^{\mu\nu}\mathbf{D}_{\mu}\mathbf{D}_{\nu}u\label{eq:cA-426}
\end{equation}

其中 $u$ 是 $M$ 上的标量函数. 或者, 在局部坐标下

\begin{equation}
\Delta_{\mathbf{M}}u(x) = \frac{1}{\sqrt{|g(x)|}}\partial_{\mu}(g^{\mu\nu}\sqrt{|g(x)|}\partial_{\nu})u(x)\label{eq:cA-427}
\end{equation}

在 Lorentz 流形的特殊情形下, Laplace--Beltrami 算子称为D'Alembert 算子, 记作 $\Box_{\mathbf{M}}$. 在任意伪 Riemann 流形上, $\Delta_{\mathbf{M}}$ 关于下面标量函数 $u$、$v$ 的标量积是对称的:

\[
(u, v)_{\mathbf{M}} = \int u(x)v(x)dv_{\mathbf{M}}
\]

事实上, 以下恒等式通过对任意两个光滑、紧支集\footnote{若流形 $\mathbf{M}$ 是紧的, 则此条件自动满足.}的函数 $u$、$v$ 分部积分容易建立:

\begin{equation}
(-\Delta u, v)_{\mathbf{M}} = \int_{\mathbf{M}}\nabla u \cdot \nabla v\, dv_{\mathbf{M}} = (u, -\Delta v)_{\mathbf{M}}\label{eq:cA-428}
\end{equation}

其中 $\nabla u\cdot\nabla v = g^{ij}\partial_{i}u\partial_{j}v$. 在 u = v 的特殊情形下, 我们推得 $(-\Delta u, v)_{\mathbf{M}} = \int_{\mathbf{M}}|\nabla u|^{2}$, 其中 $|\nabla u|^{2} = \nabla u\cdot\nabla u$. 于是, 对函数 $u \in \mathcal{C}_{0}^{\infty}(\mathbf{M})$, $-\Delta = -\Delta_{M}$ 是对称的. 若流形 $M$ 是 Riemann 的, 则它是正定的. 对 Lorentz 流形则不然:$\Box_{M}$ 不是正定的.


\section{Minkowski 空间}

\mtsubsec{基本定义}

我们记为 $R^{n+1}$ 的 $n+1$ 维 Minkowski 空间由流形 $R^{n+1}$ 连同 Lorentz 度量 $m$ 与一个特殊的坐标系 $x^{\alpha}$, $\alpha = 0, 1, \ldots, n$(称为惯性坐标系)组成, 相对于该坐标系度量具有对角形式 $\mathbf{m}_{\alpha\beta} = \text{diag}(-1, 1, \ldots, 1)$. 两个惯性坐标系之间由平移或 Lorentz 变换相联系. 我们使用关于 $m$ 及其逆 $m^{-1} = m$ 的指标升降的标准几何约定, 以及重复指标上的通常求和约定. 坐标向量场 $\frac{\partial}{\partial x^{\alpha}}$ 记作 $\partial_{\alpha}$, 任意向量场记作 $X = X^{\alpha}\partial_{\alpha}$, 其中 $X^{\alpha} = X^{\alpha}(x^{0}, \ldots, x^{n})$. 注意, 相对于 $m$ 下降指标, 得 $X_{0} = -X^{0}$, 且对一切 $i = 1, \ldots, n$ 有 $X_{i} = X^{i}$. 我们用 $D$ 表示 $R^{n+1}$ 的平坦协变导数, 即对任意一次形式 $w = \omega_{\alpha}dx^{\alpha}$ 有 $D_{\alpha}\omega_{\beta} = \partial_{\alpha}\omega_{\beta}$. 我们还把时空坐标 $x^{\alpha}$ 分为时间分量 $x^{0} = t$ 与空间分量 $x = x^{i}, \ldots, x^{n}$. 注意 $t_{0} = -t$, 且对 $i = 1, \ldots, n$ 有 $x^{i} = x_{i}$.

向量 $X$ 称为类时、类光或类空的, 分别取决于 $\mathbf{m}(X, X)$ 是 $<0$、$=0$ 还是 $>0$. 相应地, 光滑曲线 $x^{\alpha}(s)$ 称为类时、类光或类空的, 如果其切向量 $\frac{dx^{\alpha}}{ds}$ 在其每一点处是类时、类光或类空的. 因果曲线可以是类时的或类光的. 类似地, 超曲面 $u(x^{0}, \ldots, x^{n})=0$ 称为类空、类光或类时的, 如果其法向量 $N^{\alpha} = -\mathbf{m}^{\alpha\beta}\partial_{\beta}u$ 分别是类时、类光或类空的. m 在类空超曲面上诱导的度量必然正定, 即是 Riemann 的. 函数 $\mathbf{t}(x^{0}, x^{1}, \ldots, x^{n})$ 称为\newterm{time-function}{时间函数}{time function}, 如果其水平超曲面 t = t 是类空的. 在零超曲面上, 诱导度量沿法向方向是退化的, 即 $\mathbf{m}(N, N) = 0$. 特别地, 水平面 u = u 为零(即类光)的函数 $\mathbf{u} = \mathbf{u}(x^{0}, \ldots x^{n})$ 必须满足\newterm{eikonal-equation}{程函方程}{eikonal equation}

\begin{equation}
\mathbf{m}^{\alpha\beta}\partial_{\alpha}\mathbf{u}\partial_{\beta}\mathbf{u} = 0\label{eq:cA-429}
\end{equation}

方程 \eqref{eq:cA-429} 也可以写成 $D_{N}N = 0$ 的形式. 我们称 $N$ 为 $u$ 的水平超曲面的\newterm{null-geodesic-generator}{零测地线生成元}{null geodesic generator}.

因果曲线在其每一点处可以是类时的或类光的. $R^{n+1}$ 的典范时间定向由向量场 $T_{0} = \partial_{0}$ 给出. 类时向量 $X$ 称为未来定向的, 如果 $\mathbf{m}(X, T_{0}) < 0$;称为过去定向的, 如果 $\mathbf{m}(X, T_{0}) > 0$. 集合 $S$ 的\newterm{causal-future}{因果未来}{causal future} $J^{+}(S)$ 由 $R^{n+1}$ 中所有能通过未来定向因果曲线与 $S$ 相连的点组成. \newterm{causal-past}{因果过去}{causal past} $\mathcal{J}^{-}(S)$ 以同样方式定义. 于是, 对点 $p = (t, x)$, $\mathcal{J}^{+}(p) = \{(t \geq t_{0}, x)\, /\, |x - x_{0}| \leq t - t_{0}\}$. 给定光滑区域 D, 其未来集合 $\mathcal{J}^{+}(D)$ 一般可能具有不光滑的边界, 这是由于\newterm{caustic}{焦散}{caustic}引起的.

我们考虑\newterm{conservative-domain}{保守区域}{conservative domain} $\mathcal{J}^{+}(D_{1})\cap\mathcal{J}^{-}(D_{2})$, 其中 $D_{1}\subset\Sigma_{1}$、$D_{2}\subset\Sigma_{2}$ 是类空超曲面. 若 $D_{1}, D_{2}$ 都是正则的, 且其非类空边界 $\mathcal{N}_1\subset\partial(\mathcal{J}^+(D_1))\setminus D_1$ 与 $\mathcal{N}_2\subset\partial(\mathcal{J}^{-}(D_2))\setminus D_2$ 都是光滑的, 则称该区域为正则的. 在特殊情形 $D_{1} = \Sigma_{1}$、$D = D_{2}\subset\Sigma_{2}$ 下, 我们得到 $\mathcal{J}^{+}(\Sigma_{1})\cap\mathcal{J}^{-}(D)$, 称为 $D$ 关于 $\Sigma_{1}$ 的\newterm{domain-of-dependence}{依赖域}{domain of dependence}, 它由 $D\subset\Sigma_{2}$ 的因果过去中位于 $\Sigma_{1}$ 未来的所有点组成. 类似地, $\mathcal{J}^{+}(D)\cap\mathcal{J}^{-}(\Sigma_{2})$(其中 $D\subset\Sigma_{1}$)称为 $D$ 关于 $\Sigma_{2}$ 的影响依赖域. 特别有用的例子由时间函数 $\mathbf{t}$ 给出, 取 $\Sigma_1 = \{(t, x)\, /\, \mathbf{t}(t, x) = t_1\}$、$\Sigma_{2} = \{(t, x)\, /\, \mathbf{t}(t, x) = t_{1}\}$ 为两个不相交的水平超曲面, $\Sigma_{2}$ 位于 $\Sigma_{1}$ 的未来.

一对类光向量场 $L, \underline{L}$ 构成零对, 如果 $\mathbf{m}(L, \underline{L}) = -2$. 零对 $e_n = L$、$e_{n+1} = \underline{L}$ 连同满足 $\mathbf{m}(L, e_a) = \mathbf{m}(\underline{L}, e_a) = 0$ 且对一切 $a, b = 1, \ldots, n-1$ 有 $\mathbf{m}(e_a, e_b) = \delta_{ab}$ 的向量场 $e_1, \ldots, e_{n-1}$, 称为零标架. 零对

\begin{equation}
\mathbf{L} = \partial_{t} + \partial_{r}, \qquad \underline{\mathbf{L}} = \partial_{t} - \partial_{r}\label{eq:cA-430}
\end{equation}

(其中 $r = |x|$, $\partial_{r} = x^{i}/r\, \partial_{i}$)称为典范的. 类似地, 满足 $e_{n} = L$、$e_{n+1} = \underline{L}$ 的零标架 $e_{1}, \ldots, e_{n+1}$ 称为典范零标架. 此时 $e_{1}, \ldots, e_{n-1}$ 在任意点构成经过该点、$t$ 与 $r$ 恒定的球面 $S_{t, r}$ 的标准正交基. 还注意到 $L$ 是与 u = t - r 相关联的零测地线生成元, 而 $\underline{L}$ 是 $\underline{u} = t + r$ 的零测地线.

\mtsubsec{共形 Killing 向量场}

设 $x^{\mu}$ 是 Minkowski 空间 $R^{n+1}$ 的惯性坐标系. 以下是 $R^{n+1}$ 的全部等距与共形等距.

\begin{enumerate}
\renewcommand{\labelenumi}{\arabic{enumi}.}
\item \textbf{平移:}对任意给定的向量 $a = (a^{0}, a^{1}, \ldots, a^{n}) \in \mathbb{R}^{n+1}$,
\end{enumerate}

\[
x^{\mu} \rightarrow x^{\mu} + a^{\mu}
\]

\begin{enumerate}
\renewcommand{\labelenumi}{\arabic{enumi}.}
\setcounter{enumi}{1}
\item \textbf{Lorentz 旋转:}给定任意 $\Lambda = \Lambda_{\sigma}^{\rho} \in \mathbf{O}(1, n)$,
\end{enumerate}

\[
x^{\mu} \to \Lambda_{\nu}^{\mu}x^{\nu}
\]

\begin{enumerate}
\renewcommand{\labelenumi}{\arabic{enumi}.}
\setcounter{enumi}{2}
\item \textbf{伸缩:}给定任意非零实数 $\lambda$,
\end{enumerate}

\[
x^{\mu} \rightarrow \lambda x^{\mu}
\]

\begin{enumerate}
\renewcommand{\labelenumi}{\arabic{enumi}.}
\setcounter{enumi}{3}
\item \textbf{反演:}考虑变换 $x^{\mu} \rightarrow I(x^{\mu})$, 其中
\end{enumerate}

\[
I(x^{\mu}) = \frac{x^{\mu}}{(x, x)}
\]

对一切满足 $(x, x) \neq 0$ 的点 $x \in \mathbb{R}^{n+1}$ 有定义.

前两类变换是 $R^{n+1}$ 的等距, 由它们生成的群称为 Poincaré 群. 后两类变换是共形等距. 由上述所有变换生成的群称为共形群. 事实上, Liouville 定理(其无穷小版本将在后面证明)断言它就是 $R^{n+1}$ 的全部共形等距群.

下面我们列出生成上述变换的 Killing 与共形 Killing 向量场.

\begin{enumerate}
\renewcommand{\labelenumi}{\roman{enumi}.}
\item $x^{\mu}$ 方向平移的生成元, $\mu = 0, 1, \ldots, n$:
\end{enumerate}

\[
\mathbf{T}_{\mu} = \frac{\partial}{\partial x^{\mu}}
\]

\begin{enumerate}
\renewcommand{\labelenumi}{\roman{enumi}.}
\setcounter{enumi}{1}
\item $(\mu, \nu)$ 平面内 Lorentz 旋转的生成元:
\end{enumerate}

\[
\mathbf{L}_{\mu\nu} = x_{\mu}\partial_{\nu} - x_{\nu}\partial_{\mu}
\]

\begin{enumerate}
\renewcommand{\labelenumi}{\roman{enumi}.}
\setcounter{enumi}{2}
\item 伸缩变换的生成元:
\end{enumerate}

\[
\mathbf{S} = x^{\mu}\partial_{\mu}
\]

\begin{enumerate}
\renewcommand{\labelenumi}{\roman{enumi}.}
\setcounter{enumi}{3}
\item 逆平移\footnote{注意, 向量场 $K_{\mu}$ 可通过对向量场 $T_{\mu}$ 作用 $I_{*}$ 得到.}的生成元:
\end{enumerate}

\[
\mathbf{K}_{\mu} = 2x_{\mu}x^{\rho}\frac{\partial}{\partial x^{\rho}} - (x^{\rho}x_{\rho})\frac{\partial}{\partial x^{\mu}}
\]

下面我们还列出这些向量场之间的换位关系,

\begin{equation}
\begin{aligned}
&[\mathbf{L}_{\alpha\beta}, \mathbf{L}_{\gamma\delta}] = \eta_{\alpha\gamma}\mathbf{L}_{\beta\delta} - \eta_{\beta\gamma}\mathbf{L}_{\alpha\delta} + \eta_{\beta\delta}\mathbf{L}_{\alpha\gamma} - \eta_{\alpha\delta}\mathbf{L}_{\beta\gamma} \\
&[\mathbf{L}_{\alpha\beta}, \mathbf{T}_{\gamma}] = \eta_{\alpha\gamma}\mathbf{T}_{\beta} - \eta_{\beta\gamma}\mathbf{T}_{\alpha} \\
&[\mathbf{T}_{\alpha}, \mathbf{T}_{\beta}] = 0 \\
&[\mathbf{T}_{\alpha}, \mathbf{S}] = \mathbf{T}_{\alpha} \\
&[\mathbf{T}_{\alpha}, \mathbf{K}_{\beta}] = 2(\eta_{\alpha\beta}\mathbf{S} + \mathbf{L}_{\alpha\beta}) \\
&[\mathbf{L}_{\alpha\beta}, \mathbf{S}] = [\mathbf{K}_{\alpha}, \mathbf{K}_{\beta}] = 0 \\
&[\mathbf{L}_{\alpha\beta}, \mathbf{K}_{\gamma}] = \eta_{\alpha\gamma}\mathbf{K}_{\beta} - \eta_{\beta\gamma}\mathbf{K}_{\alpha}
\end{aligned}\label{eq:cA-431}
\end{equation}

记 $\mathcal{P}(1, n)$ 为由向量场 $T_{\alpha}, L_{\beta\gamma}$ 生成的 李代数, $\underline{\mathcal{K}}(1, n)$ 为由所有向量场 $T_{\alpha}, L_{\beta\gamma}, S, K_{\delta}$ 生成的 李代数, 我们陈述 Liouville 定理的如下版本:

\begin{theorem}\label{thm:cA-4.3}
以下陈述成立.

\begin{enumerate}
\renewcommand{\labelenumi}{\arabic{enumi})}
\item $\mathcal{P}(1, n)$ 是 $\mathbb{R}^{n+1}$ 中所有 Killing 向量场的 李代数.
\item 若 $n > 1$, $\underline{\mathcal{K}}(1, n)$ 是 $\mathbb{R}^{n+1}$ 中所有共形 Killing 向量场的 李代数.
\item 若 $n = 1$, $\mathbb{R}^{1+1}$ 中所有共形 Killing 向量场的集合由如下表达式给出
\end{enumerate}

\[
f(x^{0} + x^{1})(\partial_{0} + \partial_{1}) + g(x^{0} - x^{1})(\partial_{0} - \partial_{1})
\]

其中 $f$、$g$ 是单变量的任意光滑函数.
\end{theorem}

\begin{proof}
定理第 1 部分的证明作为\hyperref[prop:cA-3.4]{\propref{prop:cA-3.4}} 的特殊情形立即得出. 由 \eqref{eq:cA-417}, 因 R = 0 且 $X$ 是 Killing 的, 我们有

\[
D_{\mu}D_{\nu}X_{\lambda} = 0.
\]

因此, 存在常数 $a_{\mu\nu}, b_{\mu}$ 使得 $X^{\mu} = a_{\mu\nu}x^{\nu} + b_{\mu}$. 由于 $X$ 是 Killing 的, $D_{\mu}X_{\nu} = -D_{\nu}X_{\mu}$, 这意味着 $a_{\mu\nu} = -a_{\nu\mu}$. 从而 $X$ 可以写成向量场 $T_{\alpha}, L_{\beta\gamma}$ 的实系数线性组合.

现在设 $X$ 是共形 Killing 向量场. 存在函数 $\Omega$ 使得

\begin{equation}
{}^{(X)}\pi_{\rho\sigma} = \Omega\eta_{\rho\sigma}\label{eq:cA-432}
\end{equation}

由 \eqref{eq:cA-417} 与 \eqref{eq:cA-418} 推得

\begin{equation}
D_{\mu}D_{\nu}X_{\lambda} = \frac{1}{2}\left(\Omega_{, \mu}\eta_{\nu\lambda} + \Omega_{, \nu}\eta_{\mu\lambda} - \Omega_{, \lambda}\eta_{\nu\mu}\right)\label{eq:cA-433}
\end{equation}

对 \eqref{eq:cA-433} 两边关于 $\mu, \nu$ 取迹, 我们推得

\begin{equation}
\begin{aligned}
\Box X_{\lambda} &= -\frac{n-1}{2}\Omega_{, \lambda} \\
D^{\mu}X_{\mu} &= \frac{n+1}{2}\Omega
\end{aligned}\label{eq:cA-434}
\end{equation}

再对第一个方程作用 $D^{\lambda}$, 对第二个方程作用 $\Box$, 相减得

\begin{equation}
\Box\Omega = 0\label{eq:cA-435}
\end{equation}

对 \eqref{eq:cA-434} 的第一个方程作用 $D_{\mu}$ 并利用 \eqref{eq:cA-435}, 得

\begin{equation}
\begin{aligned}
(n-1)D_{\mu}D_{\lambda}\Omega &= \frac{n-1}{2}(D_{\mu}D_{\lambda}\Omega + D_{\lambda}D_{\mu}\Omega) = -\Box(D_{\mu}X_{\lambda} + D_{\lambda}X_{\mu}) \\
&= -(\Box\Omega)\eta_{\mu\lambda} = 0
\end{aligned}\label{eq:cA-436}
\end{equation}

因此对 $n \neq 1$, 有 $D_{\mu}D_{\lambda}\Omega = 0$. 这意味着 $\Omega$ 必为 $x^{\mu}$ 的线性函数. 于是我们可以找到常系数的线性组合 $cS + d^{\alpha}K_{\alpha}$, 使得 $X - (cS + d^{\alpha}K_{\alpha})$ 的变形张量为零. 确实如此, 因为 $^{(S)}\pi = 2\eta$ 且 $^{(K_{\mu})}\pi = 4x_{\mu}\eta$. 因此 $X - (cS + d^{\alpha}K_{\alpha})$ 是 Killing 的, 由定理的第一部分即得结论.

第 3 部分可以通过求解 \eqref{eq:cA-432} 容易导出. 事实上, 设 $X = a\partial_0 + b\partial_1$, 我们得到 $2D_0X_0 = -\Omega$、$2D_1X_1 = \Omega$、$D_0X_1 + D_1X_0 = 0$. 因此 $a$、$b$ 满足方程组

\[
\frac{\partial a}{\partial x^{0}} = \frac{\partial b}{\partial x^{1}}, \qquad \frac{\partial b}{\partial x^{0}} = \frac{\partial a}{\partial x^{1}}.
\]

于是一次形式 $adx^0 + bdx^1$ 是恰当的, $adx^0 + bdx^1 = d\phi$, 并且 $\frac{\partial^2a}{\partial x^{0^2}} = \frac{\partial^2b}{\partial x^{1^2}}$, 即 $\Box\phi = 0$. 总之

\[
X = \frac{1}{2}\left(\frac{\partial\phi}{\partial x^{0}} + \frac{\partial\phi}{\partial x^{1}}\right)(\partial_{0} + \partial_{1}) + \frac{1}{2}\left(\frac{\partial\phi}{\partial x^{0}} - \frac{\partial\phi}{\partial x^{1}}\right)(\partial_{0} - \partial_{1})
\]

这就证明了结论.
\end{proof}


\begin{remark*}
相对于典范零对表示, 有

\begin{equation}
\mathbf{T}_{0} = 2^{-1}(\mathbf{L} + \underline{\mathbf{L}}), \qquad \mathbf{S} = 2^{-1}(\underline{u}\mathbf{L} + u\underline{\mathbf{L}}), \qquad \mathbf{K}_{0} = 2^{-1}(\underline{u}^{2}\mathbf{L} + u^{2}\underline{\mathbf{L}}).\label{eq:cA-437}
\end{equation}

$T_{0} = \partial_{t}$ 与 $\mathbf{K}_{0} = (t^{2} + |x|^{2})\partial_{t} + 2tx^{i}\partial_{i}$ 都是因果的. 这使得它们在推导能量估计时很重要. 注意 $S$ 仅在 $\mathcal{J}^{+}(0) \cup \mathcal{J}^{-}(0)$ 中是因果的.
\end{remark*}

\mtsubsec{零超曲面}

零超曲面特别重要, 因为它们对应于 Minkowski 空间中波动方程或 Maxwell 方程解的传播前沿\footnote{或更一般地在 Lorentz 时空上.}. 描述零超曲面几何的最简单方法是从余维数为一的超曲面 $S_{0} \subset \Sigma_{0}$ 出发, 其中 $\Sigma_{0}$ 是 $M^{n+1}$ 的固定类空超曲面. 在 $S_{0}$ 的每一点 $p \in S_{0}$, 恰有两个与切空间 $T_{p}(S_{0})$ 正交的类光方向. 设 $L$ 表示与 $S_{0}$ 正交的光滑类光向量场, 并考虑由 $L$ 的积分曲线生成的零测地线汇\footnote{在 Minkowski 空间中, 这些实际上是直线.}. 只要这些零测地线不相交, 该线汇就构成光滑的零超曲面 N. 我们还可以通过平行移动把 $L$ 延拓到 $N$ 的所有点. 显然 $D_{L}L = 0$、$\mathbf{m}(L, L) = 0$, 此外对一切切于 $N$ 的向量 $X$ 有 $\mathbf{m}(L, X) = 0$. 还注意到 $L$ 在乘以一个只依赖于 $S_{0}$ 的共形因子的意义下唯一确定. 对一切切于 $N$ 的向量场 $X$、Y, 定义

\begin{equation}
\gamma(X, Y) = \mathbf{m}(X, Y), \qquad \chi(X, Y) = \mathbf{m}(D_{X}L, Y)\label{eq:cA-438}
\end{equation}

它们都是对称张量, 分别称为 $N$ 的\newterm{null-fundamental-form}{零基本形式}{null fundamental form}第一与第二零基本形式. 注意 $\chi$ 在乘以与 $L$ 相联系的同一个共形因子的意义下唯一确定. 显然, 对一切切于 $N$ 的 $X$ 有 $\gamma(L, X) = \chi(L, X) = 0$, 因此它们在固定的 $p\in N$ 处只依赖于与 $L_{p}$ 横截的一个固定超平面. 定义 s(称为仿射参数), 由条件 $L(s)=1$ 及在 $S_{0}$ 上 s=0 确定. 其水平面定义了 $N$ 的\newterm{geodesic-foliation}{测地叶状结构}{geodesic foliation}. 给定 $S_{0}$ 上的坐标 $w=(\omega^{a})$, $a=1, \ldots, n-1$, 我们可以用流 $x^{\mu}(s, \omega)$ 参数化 $S_{s}$ 上的点, 该流由 $\frac{dx^{\mu}}{ds}=L^{\mu}$ 定义, 其中 $x^{\mu}(0, \omega)$ 是 $S_{0}$ 上坐标为 $w$ 的点. 令

\[
\gamma_{ab} = \gamma\left(\frac{\partial}{\partial\omega^{a}}, \frac{\partial}{\partial\omega^{b}}\right), \qquad \chi_{ab} = \chi\left(\frac{\partial}{\partial\omega^{a}}, \frac{\partial}{\partial\omega^{b}}\right)
\]

表示 $\gamma$ 与 $\chi$ 相对于这些坐标的分量. 容易验证 $\frac{d}{ds}\gamma_{ab} = 2\chi_{ab}$. $S_{s}$ 的体积元由

\[
da_{S_{s}} = \sqrt{|\gamma|}\, d\omega^{1}\dots d\omega^{n-1}
\]

给出, 其中 $\gamma$ 是度量 $\gamma$ 的行列式. 注意 $\frac{d}{ds}\log|\gamma| = \gamma^{ab}\frac{d}{ds}\gamma_{ab} = 2\operatorname{tr}\chi$, 其中 $\operatorname{tr}\chi = \gamma^{ab}\chi_{ab}$ 是零超曲面的\newterm{expansion-coefficient}{膨胀系数}{expansion coefficient}. 因此,

\[
\frac{d}{ds}\sqrt{|\gamma|} = \operatorname{tr}\chi\sqrt{|\gamma|}.
\]

总体积 $|S_{s}|$ 的变化率由如下公式给出:

\begin{equation}
\frac{d}{ds}|S_{s}| = \int_{S_{s}}\operatorname{tr}\chi\, da_{S_{s}}\label{eq:cA-439}
\end{equation}

我们还注意到 $\chi$ 满足如下 Riccati 型方程,

\begin{equation}
\frac{d}{ds}\chi + \chi^{2} = 0\label{eq:cA-440}
\end{equation}

它可以显式积分. 于是可以验证, 若在 $S_{0}$ 的某点处 $\operatorname{tr}\chi(0, \omega_{0}) < 0$, 则 $\operatorname{tr}\chi(s, \omega_{0})$ 可能在有限的 $s>0$ 值处变为 $-\infty$. 这一现象对应于焦散的形成.

N 上的任意叶状结构 $S_{v}$ 可以用 $v(s, \omega)$ 参数化, 其中 $(s, \omega)$ 是上面定义的测地坐标. 我们把 $\Omega = \frac{dv}{ds}$ 称为该叶状结构的\newterm{null-lapse-function}{零时距函数}{null lapse function}, 并用 $\gamma'$ 与 $\chi'$ 表示 $\gamma, \chi$ 在 $S_{v}$ 上的限制. 若 $X$ 是切于测地叶状结构 $S_{s}$ 的向量场, 则 $X' = X - \Omega^{-1}X(v)L$ 切于 $S_{v}$. 于是, 若 $X$、$Y$ 切于 $S_{s}$, 则 $\gamma(X, Y) = \gamma(X', Y')$, $\chi(X', Y') = \chi(X, Y)$. 相对于坐标 $(v, \omega)$, 我们有

\[
\gamma_{ab}^{\prime} = \gamma_{ab}, \qquad \chi_{ab}^{\prime} = \chi_{ab}.
\]

为了在零超曲面 $N$ 上定义体积元, 我们选择任意叶状结构 v, 其零时距函数为 $\frac{dv}{ds} = \Omega$, 诱导度量为 $\gamma$, 并令

\begin{equation}
da_{\mathcal{N}} = \Omega^{-1}da_{S_{v}}dv\label{eq:cA-441}
\end{equation}

其中 $da_{S_{v}}$ 表示由 $\gamma$ 在 $S_{v}$ 上诱导的面积元. 该定义不依赖于叶状结构的选取.

\mtsubsec{能量-动量张量}

$R^{n+1}$ 中的\newterm{energy-momentum-tensor}{能量-动量张量}{energy-momentum tensor}是满足正能量条件的对称二阶张量 $Q_{\alpha\beta}$:

\[
Q(X, Y) \geq 0,
\]

对一切因果的、未来定向的 $X$、Y. 若

\begin{equation}
D^{\beta}Q_{\alpha\beta} = 0\label{eq:cA-442}
\end{equation}

我们称 $Q$ 是\newterm{divergenceless}{无散度的}{divergenceless}. 给定任意向量场 X,

\[
D^{\alpha}(Q_{\alpha\beta}X^{\beta}) = Q^{\alpha\beta}D_{\alpha}X_{\beta} = \frac{1}{2}Q^{\alpha\beta}\, {}^{(X)}\pi_{\alpha\beta},
\]

其中 ${}^{(X)}\pi = \mathcal{L}_{X}\mathbf{m}$ 表示 $X$ 的变形张量. 回忆 ${}^{(X)}\pi_{\alpha\beta} = \partial_{\alpha}X_{\beta} + \partial_{\beta}X_{\alpha}$. 在 $X$ 是 Killing 向量场的特殊情形下, 即 ${}^{(X)}\pi = 0$, 我们推得

\begin{equation}
D^{\alpha}(Q_{\alpha\beta}X^{\beta}) = 0.\label{eq:cA-443}
\end{equation}

若 $X$ 是共形 Killing 的且 $Q$ 无迹, 即 $\mathbf{m}^{\alpha\beta}Q_{\alpha\beta} = 0$, 同样的恒等式成立.

一个典型的守恒律是这样得到的:在正则保守时空区域(见第 4.1 节)$\mathcal{J}^{+}(D_{1})\cap\mathcal{J}^{-}(D_{2})$ 上积分后面的恒等式并应用 Stokes 定理, 其中光滑类空边界 $D_{i}\subset\Sigma_{i}$, 零边界 $\mathcal{N}_i$, $i = 1, 2$. 我们用 $T_{1}, T_{2}$ 表示类空超曲面 $\Sigma_1, \Sigma_2$ 的未来单位法向量, 并选取零法向量 $L_{1}, L_{2}$, 使得沿边界 $D_{i}\subset\Sigma_{i}$, $i = 1, 2$ 有 $\mathbf{m}(L_i, T_i) = -1$. 为简单起见, 只要不会引起混淆, 我们用 $T$ 表示两个类时法向量, 用 $L$ 表示两个零法向量.

\begin{proposition}\label{prop:cA-4.6}
设 $Q_{\alpha\beta}$ 是无散度的能量-动量张量, X 是上述正则保守区域 $\mathcal{J}(D_{1}, D_{2})$ 邻域中的 Killing 向量场. 则

\begin{equation}
\int_{\mathcal{N}_{2}}Q(X, L) + \int_{D_{2}}Q(X, T) = \int_{\mathcal{N}_{1}}Q(X, L) + \int_{D_{1}}Q(X, T)\label{eq:cA-444}
\end{equation}
\end{proposition}

积分分别关于零超曲面 $N_{1}, N_{2}$ 上的面积元 $da_{N}$ 以及 $m$ 在 $\Sigma_{1}, \Sigma_{2}$ 上诱导的 Riemann 度量的面积元进行. 注意, 若 $X$ 是因果的, 则所有被积函数都是正的. 若 $X$ 是共形 Killing 的且 $Q$ 无迹, 恒等式 \eqref{eq:cA-444} 仍然成立.

\begin{proof}
设 $P_{\alpha} = Q_{\alpha\beta}X^{\beta}$. 根据 \eqref{eq:cA-443} 有 $D^{\alpha}P_{\alpha} = 0$.
\end{proof}

对于依赖域 $\mathcal{J}^{+}(\Sigma_{1})\cap\mathcal{J}^{-}(D\subset\Sigma_{2})$, 或影响域 $\mathcal{J}^{+}(D\subset\Sigma_{1})\cap\mathcal{J}^{-}(\Sigma_{2})$($\Sigma_{2}$ 位于 $\Sigma_{1}$ 的未来), 结果可以简化. 我们在 $\partial D\subset\Sigma_{2}$ 上以条件 $\mathbf{m}(L, T) = -1$ 规范化 L, 其中 $T$ 表示 $\Sigma_{1}, \Sigma_{2}$ 的单位法向量.

\begin{corollary}\label{cor:cA-4.7}
若 $Q$ 无散度, X 是 Killing 的, 且 $D \subset \Sigma_2$, 则

\begin{equation}
\int_{\mathcal{N}}Q(X, L) + \int_{D \subset \Sigma_{2}}Q(X, T) = \int_{\mathcal{J}^{-}(D)\cap\Sigma_{1}}Q(X, T)\label{eq:cA-445}
\end{equation}

类似地, 若 $D\subset\Sigma_{1}$, 则

\begin{equation}
\int_{\mathcal{N}}Q(X, L) + \int_{D \subset \Sigma_{1}}Q(X, T) = \int_{\mathcal{J}^{+}(D)\cap\Sigma_{2}}Q(X, T)\label{eq:cA-446}
\end{equation}

若 $X$ 是共形 Killing 的且 $Q$ 无迹, 恒等式仍然成立.
\end{corollary}

\section*{参考文献}
\markboth{参考文献}{参考文献}
\addcontentsline{toc}{section}{参考文献}
\begin{mybibliography}{Wa}
\bibitem[Car]{car} S. M. Carroll, \emph{Lecture Note on General Relativity}, arXiv:gr-qc/9712019v1, Dec 1997.
\bibitem[E]{e} L. C. Evans, \emph{Partial Differential Equations}, Providence: American Mathematical Society, 1998.
\bibitem[Fr-Io]{fr-io} F. G. Friedlander and M. Ioshi, \emph{An introduction to the theory of distributions}, Cambridge University Press, second edition, 1998.
\bibitem[G-T]{g-t} D. Gilbarg and N. Trudinger, \emph{Elliptic partial differential equations of second order}, Springer Verlag, second edition, 1983.
\bibitem[Ge-S]{ge-s} I. M. Gelfand and G. E. Shilov, \emph{Generalized Functions, Volume I: Properties and Operations}, Academic Press, 1964.
\bibitem[Ho]{ho} L. Hörmander, \emph{The Analysis of linear partial differential operators}, vols 1-4, Springer Verlag, 1983, 1985.
\bibitem[Ha-E]{ha-e} S. W. Hawking and G. F. R. Ellis, \emph{The Large Scale Structure of Space-Time}, Cambridge University Press, 1973.
\bibitem[J]{j} F. John, \emph{Partial Differential Equations}, Providence: American Mathematical Society, fourth edition, 1982.
\bibitem[K-T]{k-t} Markus Keel and Terrence Tao, ``Endpoint Strichartz Estimates'', \emph{Am. J. Math.} 120 (1998), 955--980.
\bibitem[Kl-Rodn2]{kl-rodn2a} S. Klainerman and I. Rodnianski, Causal geometry of Einstein-vacuum spacetimes with finite curvature flux, \emph{Invent. Math.} 159 (2005), 437-529.
\bibitem[Kl-Rodn2]{kl-rodn2b} S. Klainerman and I. Rodnianski, A geometric approach to the Littlewood-Paley theory, \emph{Geom. Funct. Anal.} 16 (2006), 126--163.
\bibitem[Kl-Rodn3]{kl-rodn3} S. Klainerman and I. Rodnianski, Sharp trace theorems for null hypersurfaces on Einstein metrics with finite curvature flux, \emph{Geom. Funct. Anal.} 16 (2006), 164--229.
\bibitem[Pet]{pet} P. Petersen, \emph{Riemannian Geometry}, GTM 171, Springer, 1998.
\bibitem[Sp]{sp} M. Spivak, \emph{A Comprehensive introduction to Differential Geometry}, vols I-V, Wilmington: Publish or Perish.
\bibitem[T]{t} M. E. Taylor, \emph{Partial Differential Equations}, Springer Verlag, Texts in Applied Mathematics, vol. 23, 1996.
\bibitem[Wa]{wa} R. M. Wald, \emph{General Relativity}, The University of Chicago Press, 1984.
\end{mybibliography}
